\documentclass[11pt, letterpaper]{article}
\usepackage[margin=1in]{geometry}
\usepackage{graphicx}
\usepackage{booktabs}
\usepackage{amsmath}
\usepackage{hyperref}
\usepackage{xcolor}
\usepackage{float}
\usepackage{enumitem}

\usepackage{caption}
\usepackage[compact]{titlesec}

\hypersetup{colorlinks=true, linkcolor=blue!60!black, urlcolor=blue!60!black}
\setlength{\parskip}{5pt}
\setlength{\parindent}{0pt}
\linespread{1.15}
\titlespacing*{\section}{0pt}{12pt}{4pt}
\titlespacing*{\subsection}{0pt}{8pt}{3pt}
\captionsetup{font=small, labelfont=bf, skip=3pt}
\setlist[itemize]{nosep, leftmargin=1.2em, topsep=3pt, itemsep=2pt}

\title{\vspace{-1.5cm}\textbf{Evaluating Bitcoin Mining for\\Carbon Emission Reduction}}
\author{Keshav Krishnan\\\small ASSIP Screening Task 1 \quad|\quad March 2026}
\date{}

\begin{document}
\maketitle
\vspace{-0.4cm}

\noindent\fbox{\parbox{\dimexpr\textwidth-2\fboxsep-2\fboxrule}{%
\small\textbf{Summary:} I tested five claims from Li, Ren, and Bellos (2026) using EIA data for Texas (2019--2023) and EIA-930 hourly grid data. Four hold up: demand and supply are weakly correlated ($r = 0.19$), renewable supply is 1.5$\times$ more volatile than demand (3.3$\times$ hourly), solar+wind capacity grew 81\%, and CO$_2$ intensity fell 10.2\%. One does not: capacity factors follow a Beta distribution, not uniform (KS $p < 0.001$). But the paper's conclusions still survive. Switching from monthly to hourly data flips which outcome the model predicts. I also check how well the model fits different electricity market types.
}}

\vspace{6pt}

% ── 1. DATA ───────────────────────────────────────────
\section{Data}

The paper (Li, Ren, and Bellos, 2026; hereafter ``the paper'') argues that Bitcoin mining can reduce emissions by soaking up extra renewable electricity. To test this, I used three public EIA datasets:

\begin{itemize}
\item \textbf{EIA Form 860}, sheet ``3\_1\_Generator (Operable)'': Plant Code, State, Maximum Capacity (MW), Energy Source 1.
\href{https://www.eia.gov/electricity/data/eia860/}{eia.gov/electricity/data/eia860/}
\item \textbf{EIA Form 923}, ``Page 1 Generation and Fuel Data'': Plant Id, State, Reported Fuel Type Code, monthly Net Generation (MWh).
\href{https://www.eia.gov/electricity/data/eia923/}{eia.gov/electricity/data/eia923/}
\item \textbf{EIA-930 Hourly Grid Monitor}: Hourly demand and generation by source for ERCOT ($n = 8{,}760$ hours, 2023). \href{https://www.eia.gov/electricity/gridmonitor/}{eia.gov/electricity/gridmonitor/}
\item \textbf{EPA eGRID} (2022, Table 1): ERCOT CO$_2$ emission factor, 0.386 MT/MWh fossil average.
\href{https://www.epa.gov/egrid}{epa.gov/egrid}
\end{itemize}

I grouped fuels by EIA code: SUN/WND $\to$ solar+wind, hydro/geo/biomass $\to$ other renewable, NUC $\to$ nuclear, everything else $\to$ fossil. Focus on Texas/ERCOT (largest grid with both mining and renewables), 2019--2023.


% ── 2. ANALYSES ───────────────────────────────────────
\section{Analyses and Findings}

\subsection{Does mining coincide with more renewable capacity? (Proposition 1)}

\textbf{What I computed:} Total maximum capacity by fuel type and year from Form 860, added up for Texas.

\begin{table}[H]
\centering\small
\begin{tabular}{@{}lrrrrr@{}}
\toprule
& 2019 & 2020 & 2021 & 2022 & 2023 \\
\midrule
Solar+Wind (GW) & 30.5 & 35.0 & 43.2 & 50.7 & 55.4 \\
Fossil (GW) & 98.6 & 97.6 & 100.5 & 102.4 & 103.9 \\
SW addition rate (MW/yr) & -- & 4,480 & 8,221 & 7,498 & 4,661 \\
Renewable share of gen & 18.8\% & 21.9\% & 24.2\% & 26.5\% & 27.4\% \\
CO$_2$ intensity (kg/MWh) & 280 & 268 & 260 & 253 & 252 \\
\bottomrule
\end{tabular}
\caption{Texas capacity and generation summary (EIA 860 + 923).}
\end{table}

\textbf{Finding:} Solar+wind grew 81\% (+24.9 GW) while fossil barely changed (+5.4\%). CO$_2$ per MWh fell 10.2\% even as total generation grew 13.3\%. New additions peaked in 2021--22 ($\sim$7,860 MW/yr) then slowed to 4,661 in 2023. Growth is real but front-loaded.

I compared TX to five wind-heavy states (IL, IA, OK, KS, MN) as a quick check. TX added $\sim$15,000 MW more than expected, but it was already growing faster \textit{before} the China ban. So this is suggestive, not proof. Halaburda and Yermack (2023) find a similar pattern: TX miners using wind power shut down more often but earn more, because they buy cheap electricity during oversupply. That's the ``shock absorber'' role the paper predicts.

\begin{figure}[H]
\centering
\begin{minipage}{0.48\textwidth}
\centering
\includegraphics[width=\textwidth]{../output/fig1_texas_capacity.png}
\caption{Texas capacity by source (EIA 860).}
\end{minipage}\hfill
\begin{minipage}{0.48\textwidth}
\centering
\includegraphics[width=\textwidth]{../output/fig2_texas_generation_mix.png}
\caption{Generation mix and renewable trend.}
\end{minipage}
\end{figure}


\subsection{Are demand and renewable supply actually uncorrelated? (\S3.1)}

\textbf{What I computed:} Pearson and Spearman correlation between monthly total generation (stand-in for $\tilde{D}$) and renewable generation (stand-in for $\tilde{S} \cdot k$). 60 months.

\begin{table}[H]
\centering\small
\begin{tabular}{@{}lcc@{}}
\toprule
Test & Value & $p$-value \\
\midrule
Pearson $r$ & 0.192 & 0.142 \\
Spearman $\rho$ & 0.180 & 0.168 \\
\bottomrule
\end{tabular}
\caption{Demand-supply correlation (monthly, $n = 60$). Neither is significant at 5\%.}
\end{table}

\textbf{Finding:} Demand peaks in summer (AC). Wind peaks in spring. Different seasonal cycles mean high-demand months tell you little about renewable output. Li et al.\ argue this mismatch is what makes mining useful as a flexible load. Low correlation doesn't mean full independence, but the scatter plot (Fig.~3) shows no obvious pattern either.


\subsection{Is renewable supply more volatile? (\S3.1)}

\textbf{What I computed:} Coefficient of variation (std/mean) for monthly total demand vs.\ monthly renewable generation.

\begin{table}[H]
\centering\small
\begin{tabular}{@{}lccc@{}}
\toprule
& Mean (TWh/mo) & Std (TWh/mo) & CV \\
\midrule
Total demand & 41.9 & 6.4 & 0.153 \\
Renewable generation & 10.3 & 2.3 & 0.222 \\
\midrule
\multicolumn{3}{@{}l}{Ratio (renewable / demand)} & 1.45$\times$ \\
\bottomrule
\end{tabular}
\caption{Supply is roughly 1.5$\times$ more variable than demand (monthly).}
\end{table}

\textbf{Finding:} Renewable output swings more than demand month-to-month. This supports the need for a flexible load. The gap gets bigger when you zoom in: hourly wind CV is 0.51 vs.\ 0.16 monthly (a 3.3$\times$ ratio). Monthly data hides how much volatility miners would actually face.

\begin{figure}[H]
\centering
\begin{minipage}{0.48\textwidth}
\centering
\includegraphics[width=\textwidth]{../output/fig3_demand_supply_correlation.png}
\caption{Demand vs.\ renewable supply. Color = month.}
\end{minipage}\hfill
\begin{minipage}{0.48\textwidth}
\centering
\includegraphics[width=\textwidth]{../output/fig4_seasonal_patterns.png}
\caption{Seasonal demand and supply patterns.}
\end{minipage}
\end{figure}


\subsection{Is the uniform distribution assumption for $\tilde{S}$ valid? (\S3.6.2)}

\textbf{What I computed:} Capacity factor (CF $=$ actual output / max possible output) for each TX renewable generator. KS test against Uniform and Beta fit for comparison.

\begin{table}[H]
\centering\small
\begin{tabular}{@{}lcc@{}}
\toprule
& Wind ($n = 949$) & Solar ($n = 386$) \\
\midrule
Mean CF & 0.324 & 0.211 \\
Std & 0.118 & 0.072 \\
Range & 0.01--0.56 & 0.01--0.33 \\
\% within 1 std of mean & 69\% & 72\% \\
KS stat vs.\ Uniform & 0.204 & 0.283 \\
KS $p$-value & $< 0.001$ & $< 0.001$ \\
\midrule
Fitted Beta$(\alpha, \beta)$ & Beta(2.60, 2.25) & Beta(2.58, 1.80) \\
\bottomrule
\end{tabular}
\caption{Capacity factor distributions reject uniformity; Beta fits better.}
\end{table}

\textbf{Finding:} This is where the data pushes back on the paper.
\begin{itemize}
\item Both distributions are bell-shaped, not uniform. KS test strongly rejects ($p < 0.001$)
\item The uniform assumption overestimates how often output is near-zero or near-max
\item A Beta distribution fits better. Wind: Beta(2.60, 2.25). Solar: Beta(2.58, 1.80). Both have a single peak rather than a flat line.
\item The important thing still holds: CFs vary a lot. The paper's conclusions probably still work with a Beta, but the exact cutoffs in \S3.6.3 would shift.
\end{itemize}

\begin{figure}[H]
\centering
\includegraphics[width=0.7\textwidth]{../output/fig5_capacity_factors.png}
\caption{Capacity factor distributions vs.\ uniform benchmark (green).}
\end{figure}


\subsection{Is the residual demand actually shrinking? (Eq.\ 11--12)}

\textbf{What I computed:} CO$_2$ $=$ fossil generation $\times$ 0.386 MT/MWh (EPA eGRID). Fossil share $=$ fossil / total. Linear trend over 60 months.

\begin{table}[H]
\centering\small
\begin{tabular}{@{}ll@{}}
\toprule
Metric & Value \\
\midrule
CO$_2$ intensity change & 280 $\to$ 252 kg/MWh ($-$10.2\%) \\
Fossil share trend & $-2.0$ pp/yr ($p < 0.001$, $R^2 = 0.24$) \\
Interpretation & $\max\{D - S \cdot k, 0\}$ shrinking as $k$ grows \\
\bottomrule
\end{tabular}
\caption{Fossil share declining. This is the flip side of the renewable share trend in \S2.1 (same regression). Low $R^2$ reflects monthly noise around a clear downward trend.}
\end{table}

\textbf{Finding:} This is the most direct test of the paper's prediction, and it holds. Fossil's share is steadily shrinking as renewable capacity grows.

\begin{figure}[H]
\centering
\begin{minipage}{0.48\textwidth}
\centering
\includegraphics[width=\textwidth]{../output/fig7_texas_emissions.png}
\caption{Total CO$_2$ and emission intensity.}
\end{minipage}\hfill
\begin{minipage}{0.48\textwidth}
\centering
\includegraphics[width=\textwidth]{../output/fig8_residual_demand.png}
\caption{Residual (fossil) demand share declining.}
\end{minipage}
\end{figure}


% ── 3. CALIBRATION ────────────────────────────────────
\section{Calibration: Resolution Changes the Answer}

When I tried mapping the model's parameters to real EIA data for Texas 2023, I ran into something I didn't expect: the answer changes depending on whether you use monthly or hourly data.

\begin{table}[H]
\centering\small
\begin{tabular}{@{}lcc@{}}
\toprule
& Monthly resolution & Hourly resolution \\
\midrule
$\bar{S}$ & $\approx 0.35$ & $\approx 0.57$ \\
$C/(c\bar{S})$ & $\approx 1.04$ & $\approx 0.64$ \\
Model prediction & Case 1 (no RE investment) & Case 2 (mining helps) \\
\bottomrule
\end{tabular}
\caption{Using different time scales flips which outcome the model predicts.}
\end{table}

\begin{itemize}
\item Monthly averages compress the true range of renewable output.
\item Hourly ERCOT data (EIA-930): wind swings from near-zero to 26.3 GW within days. 1,143 hours (13\% of the year) had renewable shares above 50\%.
\item Monthly $\bar{S} \approx 0.35$ $\rightarrow$ Case 1 (no new renewables). Hourly $\bar{S} \approx 0.57$ $\rightarrow$ Case 2 (mining helps).
\item Same data, opposite conclusions. Calibration needs hourly data at minimum.
\item Hu et al.\ (2015) make the same point: how often you measure matters a lot for renewable investment conclusions.
\end{itemize}


% ── 4. REGULATORY ENVIRONMENTS ────────────────────────
\section{Model Fit Across Regulatory Environments}

The paper assumes a regulator sets prices $p$ and $P(\cdot)$. Borenstein and Bushnell (2015) lay out three structures:

\begin{itemize}
\item \textbf{Regulated} (Southeast, Midwest). Best fit. PUC sets rates and approves investments, matching Eqs.\ 14--15. Could work as a mining-specific interruptible tariff.

\item \textbf{Restructured} (ERCOT, PJM, CAISO). Prices come from markets, not a regulator. The idea still makes sense: renewables flood the grid, prices drop, miners absorb the extra. But nobody ``sets'' the price. Market power can mess up price signals (Borenstein, Bushnell, and Wolak, 2002). Capacity markets create incentives the model doesn't account for (Joskow, 2008).

\item \textbf{Municipal}. Publicly owned, trying to maximize public good. The regulator and the utility are the same entity. No misaligned incentives. Proposition 1 fits best here.
\end{itemize}

One way to think about it: let $\alpha \in [0,1]$ represent how much the utility cares about public welfare vs.\ profit. $\alpha = 1$: social planner. $\alpha = 0$: pure profit. The paper assumes something close to $\alpha = 1$.


% ── 5. LIMITATIONS ────────────────────────────────────
\section{Limitations}

\begin{itemize}
\item Monthly data misses the quick swings that cause curtailment. Hourly helps (\S3), but even that is rougher than the near-instant adjustments the paper describes.
\item EIA doesn't track mining usage separately. I can see big-picture patterns, not what miners actually do.
\item Everything here is correlation. Mining-friendly states might just be good for renewables anyway. I tried diff-in-diff to get closer to causation, but the pre-trend didn't hold. Suggestive, not proof.
\item ERCOT is unique (energy-only market, no capacity payments). One grid is one grid.
\end{itemize}


% ── REFERENCES ────────────────────────────────────────
\section{References}
\small
\begin{itemize}[nosep, leftmargin=1.2em, itemsep=3pt]
\item Borenstein, S. and J. Bushnell (2015). ``The US Electricity Industry After 20 Years of Restructuring.'' \textit{Annual Review of Economics}, 7, 437--463.
\item Borenstein, S., J. Bushnell, and F. Wolak (2002). ``Measuring Market Inefficiencies in California's Restructured Wholesale Electricity Market.'' \textit{American Economic Review}, 92(5), 1376--1405.
\item Halaburda, H. and D. Yermack (2023). ``Bitcoin Mining Meets Wall Street: A Study of Publicly Traded Crypto Mining Companies.'' NBER Working Paper No.\ 31849.
\item Hu, S., G. Souza, M. Ferguson, and W. Wang (2015). ``Capacity Investment in Renewable Energy Technology with Supply Intermittency: Data Granularity Matters!'' \textit{Manufacturing \& Service Operations Management}, 17(4), 480--494.
\item Joskow, P. (2008). ``Lessons Learned from Electricity Market Liberalization.'' \textit{The Energy Journal}, 29(S2), 9--42.
\item Li, J., H. Ren, and I. Bellos (2026). ``Bitcoin Mining (and maybe AI training) for Carbon Emission Reduction.'' Working Paper, George Mason University.
\end{itemize}
\normalsize

\vspace{6pt}
\hrule
\vspace{4pt}
{\small\noindent\textbf{Code:} \texttt{analysis.py}, \texttt{calibration.py}, \texttt{did\_analysis.py}, \texttt{ercot\_analysis.py}. See \texttt{README.md} for setup. All data from public EIA/EPA sources.}

\end{document}
